% !TEX root = ../../../main.tex
% استخدمنا انثتين من (../) لأننا نحتاج لنخرج من مجلدين حتى نصل لـ(main.tex).

\subsection{مكونات الشبكة \el{(Network Components)}}
يتألف النموذج من العناصر التالية، والتي تم ضبط بارامتراتها بناءً على الحسابات الهندسية الموضحة في الجلسة الثانية:
\begin{enumerate}
  \item \textbf{منبع الجهد \el{(Three-Phase Source)}}:
  يمثل محطة التحويل بجهد اسمي قدره $400 V$ (بين الأطوار) وتردد $50 Hz$. تم ضبط قيمة الذروة في الموديل لتكون $220\sqrt{2} V$.\\[-10pt]
  
\begin{arabicbox}{تذكر}
  يجب ضبط الإزاحة الطورية بين المنابع الثلاثة بزاوية $120^\circ$ لتكون ($0^\circ, -120^\circ, 120^\circ$) للأطوار \el{A} و \el{B} و \el{C} على التوالي، لتجنب ظهور تيارات تسلسل صفري غير مرغوبة.
\end{arabicbox}
 

  \item \textbf{ممانعة القصر للشبكة ($Z_{sc}$)}: تم تمثيلها باستخدام بلوك \el{Series RLC Branch} بقيم محسوبة لتعكس استطاعة قصر قدرها $500 MVA$، حيث بلغت المقاومة $R_{sc} = 3.184 \times 10^{-3} \Omega$ والحث $L_{sc} = 1 mH$ لضمان ظهور تشوه الجهد بوضوح عند التحليل.
  \item \textbf{ممانعة المحولة ($Z_{T}$)}:
  تم تمثيل المحولة ($1000 KVA$) بممانعة تسلسلية منسوبة للطرف الثانوي قيمتها $R_T = 2.08 \times 10^{-3} \Omega$ و $L_T = 3.05 \times 10^{-5} H$.
  \item \textbf{ممانعة الكبل \el{(Cable Impedance)}}: كبل نحاسي بطول $230 m$ ومقطع $16 mm^2$، بممانعة قدرها $R = 0.256 \Omega$ و $L = 6.05 \times 10^{-5} H$.
  \item \textbf{الحمل غير الخطي \el{(Non-linear Load)}}:
  يتكون من جسر تقويم ديودي سداسي النبضات \el{(Universal Bridge)} يغذي حملاً مستمراً \el{(RL Load)} باستطاعة $25 KW$، حيث تبلغ مقاومة الحمل $10.6 \Omega$ وتحريضه $1 mH$ .
\end{enumerate}
